\chapter{Random Walks and Markov Chains}
\label{ch:markov}

Many random systems evolve step by step, with the next step depending only on the
present state. This chapter starts with the simplest example, the random walk, and
solves the gambler's ruin problem completely. It then develops the general theory of
Markov chains: transition matrices, classification of states, hitting times, stationary
distributions, and convergence to equilibrium.

\section{The simple random walk}

Let $X_1,X_2,\dots$ be independent with $\Prob(X_i=+1)=p$ and $\Prob(X_i=-1)=q=1-p$.
The \emph{simple random walk} started at $S_0=k$ is $S_n=k+X_1+\dots+X_n$. It moves up or
down one unit per step. For $p=\tfrac12$ the walk is \emph{symmetric}.

\subsection{Gambler's ruin}

A gambler starts with $k$ units ($0\le k\le N$) and bets one unit per round, winning
with probability $p$. Play stops at $0$ (ruin) or $N$ (target). Let $h_k$ be the
probability of reaching $N$ before $0$.

\begin{theorem}[Gambler's ruin]\label{thm:ruin}
With $r=q/p$,
\[
h_k=\begin{cases}\dfrac kN,&p=\tfrac12,\\[1em]\dfrac{1-r^k}{1-r^N},&p\ne\tfrac12.\end{cases}
\]
\end{theorem}
\begin{proof}
Conditioning on the first step (law of total probability) gives, for $1\le k\le N-1$,
\[
h_k=p\,h_{k+1}+q\,h_{k-1},\qquad h_0=0,\ h_N=1 .
\]
Writing $1=p+q$ and rearranging, $p(h_{k+1}-h_k)=q(h_k-h_{k-1})$, so with
$d_k=h_k-h_{k-1}$ we have $d_{k+1}=r\,d_k$ and hence $d_k=r^{k-1}d_1$. Summing
the differences, $h_k=\sum_{j=1}^kd_j=d_1\sum_{j=0}^{k-1}r^j$. If $r=1$ this is
$kd_1$, and $h_N=1$ gives $d_1=1/N$. If $r\ne1$ it is $d_1(1-r^k)/(1-r)$, and
$h_N=1$ gives $d_1=(1-r)/(1-r^N)$.
\end{proof}

\begin{example}
A gambler with $\$90$ aims for $\$100$ in a fair game, one dollar per bet:
$h_{90}=0.9$. In a casino roulette game with $p=18/38$ ($r=10/9$) the probability of
success is $\dfrac{1-(10/9)^{90}}{1-(10/9)^{100}}\approx(9/10)^{10}\approx0.35$,
far smaller, although the game is only mildly unfavourable: the long run of bets
amplifies the small negative drift.
\end{example}

\begin{corollary}[Ruin with an infinitely rich opponent]
Starting at $k\ge1$, the probability that the walk ever reaches $0$ is $(q/p)^k$ if
$p>\tfrac12$, and $1$ if $p\le\tfrac12$.
\end{corollary}
\begin{proof}
The walk reaches $0$ before $N$ with probability $1-h_k$. As $N\to\infty$ the events
``reach $0$ before $N$'' increase to ``ever reach $0$'', so by continuity of
probability the limit of $1-h_k$ is the answer. For $p>\tfrac12$, $r<1$ and
$1-h_k\to1-(1-r^k)=r^k$. For $p<\tfrac12$, $r>1$ and
$1-h_k=\frac{r^N-r^k}{r^N-1}\to1$. For $p=\tfrac12$, $1-h_k=1-k/N\to1$.
\end{proof}

\begin{theorem}[Expected duration of the game]\label{thm:duration}
Let $D_k$ be the expected number of rounds until play stops. Then
\[
D_k=k(N-k)\ \ (p=\tfrac12),\qquad
D_k=\frac k{q-p}-\frac N{q-p}\cdot\frac{1-r^k}{1-r^N}\ \ (p\ne\tfrac12).
\]
\end{theorem}
\begin{proof}
First-step analysis gives $D_k=1+pD_{k+1}+qD_{k-1}$ with $D_0=D_N=0$ (finiteness of
$D_k$ is shown in Chapter~\ref{ch:mart}). For $p=\tfrac12$ try $D_k=k(N-k)$:
$1+\frac12(k+1)(N-k-1)+\frac12(k-1)(N-k+1)=1+k(N-k)-1=k(N-k)$. The solution of the
recurrence with the boundary conditions is unique (the difference of two solutions
satisfies the homogeneous equation with zero boundary values, hence is $0$ by the
maximum principle), so this is it. For $p\ne\tfrac12$, a particular solution is
$k/(q-p)$ (check: $1+\frac{p(k+1)+q(k-1)}{q-p}=1+\frac{k-(q-p)}{q-p}=\frac k{q-p}$), and the
homogeneous solutions are $A+Br^k$. The boundary conditions give
$D_k=\frac k{q-p}-\frac N{q-p}\cdot\frac{1-r^k}{1-r^N}$.
\end{proof}

For the fair game, $D_k=k(N-k)$: a fair game between a gambler with $1$ unit and a
casino with $N-1$ units lasts on average $N-1$ rounds, yet the gambler is ruined with
probability $1-1/N$.

\subsection{Reflection principle and the maximum}

Visualise the walk as a path $(0,S_0),(1,S_1),\dots,(n,S_n)$ in the plane.

\begin{theorem}[Reflection principle]\label{thm:reflection}
Let $x,y>0$. The number of paths from $(0,x)$ to $(n,y)$ that touch or cross the level $0$
equals the number of paths from $(0,-x)$ to $(n,y)$.
\end{theorem}
\begin{proof}
Given a path from $(0,x)$ to $(n,y)$ that touches $0$, let $t$ be the first time it
does so, and reflect the portion before $t$ in the horizontal axis (replace $S_j$ by
$-S_j$ for $j\le t$). This yields a path from $(0,-x)$ that reaches $0$ at time $t$ and
then continues to $(n,y)$. The map is a bijection onto paths from $(0,-x)$ to $(n,y)$,
since any such path must cross $0$ (it starts below and ends above), and reflecting its
initial segment up to the first touch gives the inverse.
\end{proof}

\begin{theorem}[Distribution of the maximum]\label{thm:maxwalk}
For the symmetric simple random walk from $0$ and an integer $a\ge1$,
\[
\Prob\Bigl(\max_{0\le k\le n}S_k\ge a\Bigr)=\Prob(S_n\ge a)+\Prob(S_n>a)
=2\,\Prob(S_n>a)+\Prob(S_n=a).
\]
\end{theorem}
\begin{proof}
Split according to the final value. If $S_n\ge a$ the maximum is at least $a$ (the path
ends above or at $a$). If $S_n=b<a$ and the maximum is at least $a$, the path touches the
level $a$; by the reflection principle (reflecting across $a$, shifting the origin) the
number of such paths equals the number of paths ending at $2a-b>a$. Since all paths
have probability $2^{-n}$, $\Prob(\max\ge a,S_n=b)=\Prob(S_n=2a-b)$ for $b<a$.
Summing over $b<a$ gives $\Prob(S_n>a)$. Adding $\Prob(S_n\ge a)$ gives the claim.
\end{proof}

\begin{theorem}[Returns to the origin]\label{thm:returns}
For the symmetric walk, $\Prob(S_{2n}=0)=\binom{2n}n4^{-n}\sim1/\sqrt{\pi n}$, and the walk
returns to $0$ with probability $1$ (\emph{recurrence}).
\end{theorem}
\begin{proof}
The first claim: among the $2^{2n}$ paths, those with $n$ ups and $n$ downs number
$\binom{2n}n$; the asymptotics follow from Stirling's formula
$n!\sim\sqrt{2\pi n}(n/e)^n$ (cited). For recurrence let $u_{2n}=\Prob(S_{2n}=0)$ and
let $f_{2m}$ be the probability that the first return is at time $2m$. Decomposing on the
time of the first return, $u_{2n}=\sum_{m=1}^nf_{2m}u_{2n-2m}$ for $n\ge1$, with
$u_0=1$. In generating functions $U(s)=\sum_{n\ge0}u_{2n}s^{2n}$ and
$F(s)=\sum_{m\ge1}f_{2m}s^{2m}$ this reads $U=1+FU$, i.e.\ $F=1-1/U$. By the binomial
series, $\binom{2n}n4^{-n}=(-1)^n\binom{-1/2}n$, so $U(s)=(1-s^2)^{-1/2}$ for $\abs s<1$.
Hence $F(s)=1-\sqrt{1-s^2}\to1$ as $s\uparrow1$, i.e.\ $\sum_mf_{2m}=1$: the walk
returns with probability $1$.
\end{proof}

\begin{nuance}{recurrent but with infinite expected return time}
$F(s)=1-\sqrt{1-s^2}$ has an infinite derivative at $s=1$, so the expected time to
return to $0$ is $\infty$. The symmetric walk returns for sure, but typically after a
very long time ($\Prob(\text{first return}>2n)=\Prob(S_{2n}=0)\sim1/\sqrt{\pi n}$,
which is not summable). Such a state is called \emph{null recurrent}.
\end{nuance}

\section{Markov chains}

\begin{definition}[Markov chain]
A sequence $X_0,X_1,\dots$ of random variables taking values in a finite or countable
\emph{state space} $\mathcal S$ is a \emph{(time-homogeneous) Markov chain} if for all
$n$ and all states,
\[
\Prob(X_{n+1}=j\given X_n=i,X_{n-1}=i_{n-1},\dots,X_0=i_0)=\Prob(X_{n+1}=j\given X_n=i)=p_{ij}.
\]
The matrix $P=(p_{ij})$ is the \emph{transition matrix}: $p_{ij}\ge0$ and
$\sum_jp_{ij}=1$ for every $i$.
\end{definition}

The defining property, the \emph{Markov property}, says that given the present state,
the past is irrelevant to the future.

\begin{proposition}[Joint law and Chapman--Kolmogorov]\label{prop:ck}
With initial distribution $\mu_i=\Prob(X_0=i)$,
\[
\Prob(X_0=i_0,X_1=i_1,\dots,X_n=i_n)=\mu_{i_0}p_{i_0i_1}\cdots p_{i_{n-1}i_n}.
\]
Consequently the $n$-step transition probabilities
$p_{ij}^{(n)}=\Prob(X_n=j\given X_0=i)$ are the entries of the matrix power $P^n$, and
$p^{(m+n)}_{ij}=\sum_kp^{(m)}_{ik}p^{(n)}_{kj}$. The law of $X_n$ is the row vector
$\mu P^n$.
\end{proposition}
\begin{proof}
The product formula follows from the chain rule (Theorem~\ref{thm:chain}) and the
Markov property. Summing it over $i_1,\dots,i_{n-1}$ gives $p_{ij}^{(n)}$ as the $(i,j)$
entry of the product of $n$ copies of $P$ (matrix multiplication is exactly this sum).
The identity $P^{m+n}=P^mP^n$ is the Chapman--Kolmogorov equation. The last claim is
total probability: $\Prob(X_n=j)=\sum_i\mu_ip^{(n)}_{ij}$.
\end{proof}

\begin{example}[Two-state chain]\label{ex:twostate}
States $\{1,2\}$, $P=\begin{pmatrix}1-a&a\\b&1-b\end{pmatrix}$ with $a,b\in(0,1)$.
Solving $p^{(n+1)}_{11}=p^{(n)}_{11}(1-a)+p^{(n)}_{12}b=p^{(n)}_{11}(1-a-b)+b$ gives
\[
p^{(n)}_{11}=\frac b{a+b}+\frac a{a+b}(1-a-b)^n,
\]
which tends to $b/(a+b)$ geometrically fast. The vector $\pi=(\frac b{a+b},\frac a{a+b})$
satisfies $\pi P=\pi$.
\end{example}

\begin{example}[Further chains]
\begin{itemize}
\item \emph{Random walk on a graph}: from vertex $i$ move to a uniformly chosen
neighbour: $p_{ij}=1/\deg(i)$ if $i\sim j$.
\item \emph{Gambler's ruin chain}: states $0,\dots,N$; $0$ and $N$ absorbing.
\item \emph{Ehrenfest urn}: $N$ balls in two urns; each step a uniformly random ball is
moved to the other urn; the state is the number of balls in urn~1.
\item \emph{Card shuffling}: the state is a permutation of the deck; each shuffle
applies a random permutation.
\end{itemize}
\end{example}

\subsection{Classification of states}

\begin{definition}
State $j$ is \emph{accessible} from $i$ ($i\to j$) if $p^{(n)}_{ij}>0$ for some
$n\ge0$; $i$ and $j$ \emph{communicate} ($i\leftrightarrow j$) if each is accessible
from the other. Communication is an equivalence relation; its classes are the
\emph{communicating classes}. The chain is \emph{irreducible} if there is one class.
The \emph{period} of $i$ is $\gcd\{n\ge1:p^{(n)}_{ii}>0\}$; $i$ is \emph{aperiodic}
if the period is $1$. (Period is constant on a class.)
\end{definition}

Let $T_i=\min\{n\ge1:X_n=i\}$ be the first return time to $i$, and
$f_i=\Prob_i(T_i<\infty)$ (subscript $i$: the chain starts at $i$).

\begin{definition}[Recurrence]
State $i$ is \emph{recurrent} if $f_i=1$ and \emph{transient} if $f_i<1$. A recurrent
state is \emph{positive recurrent} if $\Exp_i[T_i]<\infty$ and \emph{null recurrent}
otherwise.
\end{definition}

\begin{theorem}[Recurrence criterion]\label{thm:reccrit}
Let $V_i$ be the total number of visits to $i$ (counting time $0$) when starting at
$i$. Then $V_i$ is geometric: $\Prob_i(V_i=m)=f_i^{m-1}(1-f_i)$. Hence
\[
i\text{ is recurrent}\iff\Prob_i(V_i=\infty)=1\iff\sum_{n=0}^\infty p^{(n)}_{ii}=\infty,
\]
and $i$ is transient iff $\sum_np^{(n)}_{ii}<\infty$.
\end{theorem}
\begin{proof}
By the Markov property, each time the chain is at $i$ it returns with probability
$f_i$, independently of the past. So $\Prob_i(V_i\ge m)=f_i^{m-1}$. If $f_i<1$ then
$V_i$ is a.s.\ finite with mean $\Exp_i[V_i]=1/(1-f_i)<\infty$; if $f_i=1$ then
$V_i=\infty$ a.s. Finally $\Exp_i[V_i]=\sum_n\Exp_i[\indic_{\{X_n=i\}}]=\sum_np^{(n)}_{ii}$.
\end{proof}

\begin{corollary}
Recurrence is a class property: if $i\leftrightarrow j$ and $i$ is recurrent, so is $j$.
In a finite chain every recurrent class is positive recurrent, and every state in a class
that is not closed (can be left) is transient. A finite irreducible chain is positive
recurrent.
\end{corollary}
\begin{proof}[Proof of the first assertion]
Choose $r,s$ with $p^{(r)}_{ij}>0$, $p^{(s)}_{ji}>0$. Then
$p^{(r+n+s)}_{jj}\ge p^{(s)}_{ji}p^{(n)}_{ii}p^{(r)}_{ij}$, and summing over $n$ shows
$\sum_np^{(n)}_{jj}=\infty$. (The other assertions follow from the criterion and the
geometric trials argument; see the exercises.)
\end{proof}

\begin{example}[Random walks in higher dimensions]
The symmetric simple random walk on $\Z^d$ is recurrent for $d=1,2$ and transient for
$d\ge3$ (P\'olya). For $d=1$ we showed this directly. For $d=3$, $\Prob(S_{2n}=0)\asymp
n^{-3/2}$, which is summable, so by Theorem~\ref{thm:reccrit} the origin is transient.
In the memorable phrase, ``a drunk man will find his way home, but a drunk bird may
get lost forever''.
\end{example}

\subsection{Hitting probabilities and times}

Let $A\subseteq\mathcal S$ and define $h_i=\Prob_i(\text{chain ever hits }A)$ and
$k_i=\Exp_i[\text{time to hit }A]$.

\begin{theorem}[Linear equations for hitting]\label{thm:hitting}
The vector $h$ satisfies $h_i=1$ for $i\in A$ and $h_i=\sum_jp_{ij}h_j$ for $i\notin A$;
it is the \emph{minimal non-negative} solution of these equations. The vector $k$ is
the minimal non-negative solution of $k_i=0$ for $i\in A$, $k_i=1+\sum_jp_{ij}k_j$ for
$i\notin A$.
\end{theorem}
\begin{proof}
That $h$ solves the equations is first-step analysis. For minimality, let $x$ be any
non-negative solution and let $h_i^{(n)}=\Prob_i(\text{hit }A\text{ within }n\text{ steps})$.
Then $h^{(0)}_i=\indic_{A}(i)\le x_i$, and if $h^{(n)}\le x$ then
$h^{(n+1)}_i=\sum_jp_{ij}h^{(n)}_j\le\sum_jp_{ij}x_j=x_i$ for $i\notin A$. So $h^{(n)}\le x$
for all $n$, and $h^{(n)}\uparrow h$. The argument for $k$ is the same.
\end{proof}

\begin{nuance}{minimality matters on infinite state spaces}
For gambler's ruin on $\{0,1,2,\dots\}$ with $p>\tfrac12$, both $h_k\equiv1$ and
$h_k=(q/p)^k$ solve the equations for hitting $0$, but the true probability is the
smaller one. In finite chains where $A$ is reached with probability $1$ from every state,
the solution is unique.
\end{nuance}

\begin{theorem}[Absorbing chains]\label{thm:absorbing}
Suppose the states are split into transient states $\mathcal T$ and absorbing states
$\mathcal A$, with $P=\begin{pmatrix}Q&R\\0&I\end{pmatrix}$ ($Q$ for transitions among
transient states, $R$ for transient$\to$absorbing). Then $I-Q$ is invertible, the matrix
$N=(I-Q)^{-1}=\sum_{n\ge0}Q^n$ has entries $N_{ij}=\Exp_i[\text{visits to }j]$, the expected
time to absorption from $i$ is $\sum_jN_{ij}$, and the absorption probabilities are
$B=NR$.
\end{theorem}
\begin{proof}
$Q^n_{ij}=\Prob_i(X_n=j)$ for transient $j$, so $\sum_nQ^n_{ij}$ is the expected number
of visits to $j$. Absorption occurs eventually with probability one (finite chain, every
transient state leads to absorption), so $Q^n\to0$, which implies the Neumann series
$\sum Q^n$ converges and equals $(I-Q)^{-1}$ (multiply $(I-Q)\sum_{m\le n}Q^m=I-Q^{n+1}$
and let $n\to\infty$). The probability of absorbing at $a$ from $i$ is
$\sum_n\sum_jQ^n_{ij}R_{ja}=(NR)_{ia}$, summing over the time of the last transient step.
\end{proof}

\section{Stationary distributions and convergence}

\begin{definition}[Stationary distribution]
A probability vector $\pi$ (row, non-negative, $\sum\pi_i=1$) is \emph{stationary} (or
invariant) if $\pi P=\pi$. If $X_0\sim\pi$ then $X_n\sim\pi$ for all $n$.
\end{definition}

\begin{theorem}[Existence and formula]\label{thm:stationary}
Let the chain be irreducible and recurrent, fix a state $i$, and define
$\gamma_j=\Exp_i\Bigl[\sum_{n=0}^{T_i-1}\indic_{\{X_n=j\}}\Bigr]$, the expected number of
visits to $j$ in one excursion from $i$. Then $\gamma P=\gamma$, $\gamma_i=1$, and
$\sum_j\gamma_j=\Exp_i[T_i]$. Consequently, if the chain is positive recurrent,
$\pi_j=\gamma_j/\Exp_i[T_i]$ is a stationary distribution and in particular
$\pi_i=1/\Exp_i[T_i]$.
\end{theorem}
\begin{proof}
Write $\gamma_j=\sum_{n\ge0}\Prob_i(X_n=j,T_i>n)$. Then
\[
(\gamma P)_k=\sum_{n\ge0}\sum_j\Prob_i(X_n=j,T_i>n)p_{jk}=\sum_{n\ge0}\Prob_i(X_{n+1}=k,T_i>n),
\]
by the Markov property (the event $\{T_i>n\}$ depends on $X_0,\dots,X_n$ only). If
$k\ne i$, then on $\{X_{n+1}=k\}$ we have $T_i>n\iff T_i>n+1$, so
$(\gamma P)_k=\sum_{n\ge0}\Prob_i(X_{n+1}=k,T_i>n+1)=\sum_{m\ge1}\Prob_i(X_m=k,T_i>m)=\gamma_k$
(the term $m=0$ is $0$ since $X_0=i\ne k$). If $k=i$,
$(\gamma P)_i=\sum_{n\ge0}\Prob_i(X_{n+1}=i,T_i>n)=\sum_{n\ge0}\Prob_i(T_i=n+1)=1=\gamma_i$,
using recurrence. Finally $\sum_j\gamma_j=\Exp_i\sum_{n<T_i}1=\Exp_i[T_i]$.
\end{proof}

\begin{theorem}[Convergence to equilibrium]\label{thm:convergence}
Let the chain be finite, irreducible and aperiodic. Then it has a unique stationary
distribution $\pi$, all $\pi_j>0$, and for some $C<\infty$ and $\rho\in(0,1)$,
\[
\max_{i}\sum_j\bigl\lvert p^{(n)}_{ij}-\pi_j\bigr\rvert\le C\rho^n\qquad\text{for all }n.
\]
\end{theorem}
\begin{proof}
\emph{Positivity of a power.} For a finite, irreducible, aperiodic chain there is $m$ such
that $p^{(m)}_{ij}>0$ for \emph{all} $i,j$. (Aperiodicity makes the set of return times
to $i$ closed under addition with gcd $1$, hence containing all large integers, a
standard number-theoretic fact; irreducibility then connects different states.) Let
$\epsilon=\min_{i,j}p^{(m)}_{ij}>0$. A stationary $\pi$ exists by Theorem~\ref{thm:stationary}
(a finite irreducible chain is positive recurrent).

\emph{Coupling.} Let $(X_n)$ start at $i$ and $(Y_n)$ start from $\pi$, independent, both
with transition matrix $P$. Let $\tau=\min\{n:X_n=Y_n\}$. In each block of $m$ steps, whatever
the states $x,y$ at the start of the block, $\Prob(X_{(k+1)m}=Y_{(k+1)m}\given\text{past})
\ge\sum_jp^{(m)}_{xj}p^{(m)}_{yj}\ge\epsilon\sum_jp^{(m)}_{xj}=\epsilon$. Hence
$\Prob(\tau>km)\le(1-\epsilon)^k$. Define $\tilde X_n=X_n$ for $n<\tau$ and $\tilde X_n=Y_n$
for $n\ge\tau$; by the strong Markov property $(\tilde X_n)$ is also a Markov chain with
matrix $P$ started at $i$ (we use this without detailed proof). Since $Y_n\sim\pi$ and $\tilde X_n$ has the same law as $X_n$,
\[
\bigl\lvert p^{(n)}_{ij}-\pi_j\bigr\rvert=\bigl\lvert\Prob(\tilde X_n=j)-\Prob(Y_n=j)\bigr\rvert
\le\Prob(\tilde X_n=j,Y_n\ne j)+\Prob(Y_n=j,\tilde X_n\ne j).
\]
Summing over $j$, each outcome with $\tilde X_n\ne Y_n$ is counted at most twice, and
$\{\tilde X_n\ne Y_n\}=\{\tau>n\}$, so
$\sum_j\abs{p^{(n)}_{ij}-\pi_j}\le2\Prob(\tau>n)\le2(1-\epsilon)^{\lfloor n/m\rfloor}$. This gives the bound with
$\rho=(1-\epsilon)^{1/m}$. The same argument applied to a second stationary $\pi'$ shows
$\pi'=\pi$ (it is stationary, so the distribution stays $\pi'$ while converging to $\pi$).
Positivity of $\pi$: $\pi_j=\sum_i\pi_ip^{(m)}_{ij}\ge\epsilon>0$.
\end{proof}

\begin{nuance}{aperiodicity is needed for convergence of $P^n$, not for existence of $\pi$}
Take $P=\begin{pmatrix}0&1\\1&0\end{pmatrix}$ (period $2$). Then $\pi=(\frac12,\frac12)$ is
stationary, but $P^n$ alternates between $P$ and $I$ and does not converge. For periodic
chains one has convergence of time averages (the ergodic theorem: the fraction of time spent
in state $j$ tends to $\pi_j$ almost surely, cited) but not of the distribution.
\end{nuance}

\begin{theorem}[Detailed balance]\label{thm:detailed}
If a probability vector $\pi$ satisfies $\pi_ip_{ij}=\pi_jp_{ji}$ for all $i,j$, then $\pi$
is stationary.
\end{theorem}
\begin{proof}
$(\pi P)_j=\sum_i\pi_ip_{ij}=\sum_i\pi_jp_{ji}=\pi_j\sum_ip_{ji}=\pi_j$.
\end{proof}

\begin{example}[Random walk on a graph]
For the walk on a finite connected graph, take $\pi_i=\deg(i)/(2\abs E)$. Detailed balance
holds: for adjacent $i,j$, $\pi_ip_{ij}=\frac{\deg i}{2\abs E}\cdot\frac1{\deg i}
=\frac1{2\abs E}=\pi_jp_{ji}$. So the stationary probability is proportional to
degree; and $\Exp_i[T_i]=2\abs E/\deg(i)$.
\end{example}

\begin{example}[Ehrenfest urn]
With $N$ balls, the stationary distribution is $\mathrm{Bin}(N,\tfrac12)$: check detailed
balance with $\pi_k=\binom Nk2^{-N}$ and $p_{k,k+1}=(N-k)/N$, $p_{k+1,k}=(k+1)/N$:
$\binom Nk\frac{N-k}N=\binom N{k+1}\frac{k+1}N$ (both equal $\frac{N!}{k!(N-k-1)!\,N}$).
\end{example}

\begin{interview}{what to do with a Markov chain problem}
(1) Define the state space so that the future depends only on the state. (2) For
expected times or probabilities of events, solve the first-step equations
(Theorem~\ref{thm:hitting}). (3) For long-run behaviour, solve $\pi P=\pi$, trying
detailed balance first. (4) For absorbing problems, compute $N=(I-Q)^{-1}$.
Typical interview examples: expected time to see a pattern, expected rolls to see all faces,
random walk on a cube, ``drunkard on a circle''.
\end{interview}

\section*{Exercises}
\addcontentsline{toc}{section}{Exercises}

\begin{enumerate}
\item A random walk on $\{0,1,2,3,4\}$ with absorbing ends $0,4$ and
fair steps. Find the absorption probabilities and expected duration from each state,
using Theorem~\ref{thm:absorbing}.
\item For the symmetric walk, show the probability that it never returns to $0$ in
the first $2n$ steps is $\Prob(S_{2n}=0)$.
\item A rat walks on the corners of a cube, moving to a uniformly random adjacent
corner each step. Find the expected time to return to the starting corner, and the expected
time to reach the opposite corner.
\item Show that for a finite irreducible chain $\Exp_j[T_i]<\infty$ for all $i,j$
(hint: from every state $i$ reachable within $m$ steps with probability $\ge\epsilon$;
geometric trials).
\item Prove that in a finite chain every closed communicating class is positive recurrent,
and every communicating class that is not closed consists of transient states.
\item Draw the transition graph and find the stationary distribution of the chain
on $\{1,2,3\}$ with $P=\begin{pmatrix}\frac12&\frac12&0\\ \frac13&\frac13&\frac13\\0&\frac12&\frac12\end{pmatrix}$.
\item Show that if $\pi$ is stationary for an irreducible chain then $\pi_i>0$ for all $i$.
\item Show that a time-reversed stationary Markov chain, $Y_n=X_{N-n}$, is Markov with
$\hat p_{ij}=\pi_jp_{ji}/\pi_i$, and that the chain is reversible iff $\hat P=P$.
\item (Challenge) Derive the expected number of rolls of a fair die to see all six faces via
the Markov chain on the number of faces seen, and compare with the coupon collector answer
$6H_6=14.7$.
\end{enumerate}
